% =========================================================
% mechlib/core/colors.tex
% ZENTRALE FARBPALETTE
% =========================================================
% Alle bibliotheksweiten Farben werden ausschliesslich hier definiert. Dadurch
% kann der visuelle Charakter der gesamten mechlib an einer Stelle angepasst
% werden, ohne einzelne Elemente zu veraendern.
%
% Konvention:
%   mechBodyColor    Fuellung allgemeiner Koerper und Fachwerkgelenke
%   mechSupportColor Fuellung von Lagerkoerpern
%   mechLoadColor    Kraefte, Momente und Streckenlasten
%   mechAxisColor    Achsen und Koordinatensysteme
%   mechRodColor     historische Stabfarbe; aktuelle Fachwerkstaebe sind schwarz
%
% Die Namen sind Teil der internen visuellen API und sollten bevorzugt statt
% harter RGB-Werte in den Elementdateien verwendet werden.
% =========================================================
% -------------------------------------------------------------------------
% Farbdefinitionen
% -------------------------------------------------------------------------
% RGB wird bewusst explizit verwendet, damit die Darstellung unabhaengig von
% Dokumentklassen und externen Farbmodellen reproduzierbar bleibt.
\definecolor{mechBodyColor}{RGB}{180,180,180}
\definecolor{mechSupportColor}{RGB}{205,205,205}
\definecolor{mechLoadColor}{RGB}{178,25,25}
\definecolor{mechAxisColor}{RGB}{6,13,141}
\definecolor{mechRodColor}{RGB}{130,130,130}
